% ===== §5 Value, Perception, and Circulation =====
\section{Value, Perception, and Circulation}
\label{sec:value}

\noindent
The two preceding sections established that perception is generated under a prior frame and that its symbolisation may generate or appropriate. Neither has yet asked what is generated, in the register of value, when perception goes well, nor what is stolen when it goes badly, nor why the everyday should be the place where the answer matters most. This section supplies the value theory the paper's argument has so far leaned on without stating. It does three things in sequence. It sets out how the perception of value, of the ethical, of the just, and of the beautiful is generated across the three registers this series has used throughout. It shows, with the political economy of the senses, that the frame under which perception is generated is itself historically produced, so that perceptual capacity is never a natural endowment but a social product. And it isolates what is peculiar to everyday aesthetic value, the low-amplitude, high-frequency, non-accumulable value that is, for exactly those reasons, the hardest to steal, before closing on the circulation of that value and the justice of its return.

\subsection{The perception of value across the registers}
\label{sub:val-registers}

Perception does not merely deliver objects; it delivers them as valuable, as fitting or unfitting, as beautiful or mean, and the different kinds of value it delivers are generated differently. The three-register account this series has developed can be brought to bear directly on perception, and doing so clarifies what an education of perception would be educating. In the register of the imaginary, value is perceived as an idealised completeness projected onto the object, the beauty that promises to fill a lack and to make the perceiver whole; this is the value of the perfect image, and it is exhausting in the precise sense that it can never be possessed without disappointing, since the completeness it promises was the perceiver's own projected wish. In the register of the symbolic, value is perceived as a place in a system of differences, the worth a thing has by its rank, its price, its position in an order of comparison; this is the value that can be counted and exchanged, and it is, as the previous section found, the value most readily appropriated, because what is commensurable can be transferred. In the register of the real, value is perceived around the empty place, as the incommensurable worth of what cannot be ranked, priced, or substituted, the value that gathers around the \zh{das Ding} at the centre of a genuine attachment and that no symbolic equivalent can discharge. The perception of the beautiful, of the ethical, and of the just each runs across these registers. The beautiful may be perceived as imaginary completeness, as symbolic distinction, or as real and incommensurable singularity; the ethical may be perceived as an idealised purity, as conformity to a coded rule, or as the irreducible claim of a particular other; the just may be perceived as an imagined harmony, as an accountable balance of exchangeable goods, or as the real and unsettleable demand that no participant be reduced to fuel. An education of perception is, among other things, an education in which register one's perception of value chiefly moves in, and the wager of this paper is that a cultivated everyday perception is one drawn increasingly toward the register of the real, where alone, on this series' account, endogenously good value is carried.

\subsection{The political economy of the senses}
\label{sub:val-politicaleconomy}

The frame under which perception is generated is not given by nature, and the claim that it is plastic, which the next spine will establish physiologically, has a social form that must be stated first, because it bears directly on justice. Marx stated it in the 1844 manuscripts with a directness that has not been improved upon: the forming of the five senses, he wrote, is a labour of the entire history of the world down to the present \citep{marx1844}. The senses are not a fixed biological endowment through which a ready-made world is received; they are themselves produced, historically and socially, so that what a perceiver is capable of perceiving is the sediment of a history. Marx drew the consequences that matter here. The perceiver bound by crude practical need has only a restricted sense: to the starving, food exists only in its abstract character as food and not in its human form; the care-burdened and poverty-stricken have no sense for the finest play; the dealer in minerals sees the commercial value of the stone but not its beauty, having, in Marx's phrase, no mineralogical sense. Perceptual incapacity is here shown to be produced by a mode of life, and aesthetic perception in particular to be conditioned by the material circumstances that free or bind the senses. This has a contemporary extension that the argument requires, since the mode of production under which perception is now formed has made a commodity of perception's own scarce condition. Attention, Simon observed, is what a wealth of information consumes, so that an information-rich world is an attention-poor one, and the efficient allocation of a now-scarce attention becomes an economic problem \citep{simon1971}. The everyday perception this paper cares about is conducted under precisely this regime: an economy organised to capture and resell attention systematically draws off the perceptual capacity that everyday aesthetic life requires, outsourcing the perceiver's seeing to a flow of engineered images and leaving the light on the wall unattended not by accident but by design. The political economy of the senses therefore joins the diagnosis of \xref{sec:diagnostics} to the theory of justice the paper will reach: the habituation and numbing named there are not merely private failures of attention but the subjective face of a social production of perceptual incapacity, and the appropriative symbolisation of \xref{sec:symbolic} is the cultural form through which that production operates. Who can perceive the everyday, and how well, is in part a question of political economy.

\subsection{The value of everyday aesthetic perception}
\label{sub:val-everyday}

Everyday aesthetic value has a character of its own, and naming it precisely is one of this paper's contributions, because that character is what makes everyday aesthetic value at once endangered and peculiarly resistant to theft. It is, first, low in amplitude. No single instance of it makes a large claim; the value of this morning's light is small, and makes no demand comparable to that of a masterwork. It is, second, high in frequency. What is small in each instance is nonetheless delivered continuously, so that the aggregate everyday aesthetic value of a life is very large, carried in a multitude of minor increments rather than a few major ones. It is, third, non-accumulable. Everyday aesthetic value cannot be stored, collected, or held as a stock; the light perceived this morning does not deposit a quantity of value that can be drawn on tomorrow, for its value was in the perceiving and passed with it, and tomorrow's must be generated anew. These three features together yield a fourth, which is decisive for the argument. Everyday aesthetic value is, of all aesthetic value, the hardest to appropriate. What is low in amplitude offers little worth expropriating in any single instance; what is non-accumulable cannot be stored and therefore cannot be stockpiled, converted, or transferred; and what must be generated afresh at each occurrence cannot be captured once and possessed. The masterwork can be bought, insured, and hung as a sign of its owner's standing, its symbolic value abstracted from any perceiving and made to circulate as capital. The light on the wall cannot. Its value exists only in the recurrent generative event of its being perceived, resists conversion into a possessable and exchangeable token, and so remains, structurally, close to the register of the real. This is why the everyday, for all its exposure to habituation, is the privileged site of endogenously good aesthetic value. Precisely because it is too small, too frequent, and too fleeting to be worth stealing, it is the aesthetic value least corrupted by appropriation, and the one whose generation an education of perception can most hope to protect.

\subsection{Circulation and the justice of return}
\label{sub:val-circulation}

Value that is generated must circulate, and the character of its circulation is the difference between a good and a bad aesthetic life. This series has drawn the distinction sharply, and it applies without alteration here. In a good cycle, the value generated in perception returns to the perceivers with an endogenous surplus produced by the internal tension of their shared attending, a surplus possessed by no one and settling nothing, so that the cycle can turn again with the real value at its centre still unfilled. In a bad cycle, the value either fails to return, dissipated into habituation, or returns only by external extraction, one party's perceptual life drawn off to feed another's. The just cycle, in the terms developed with Eglash, is the one in which every party to the generative process is a beneficiary of it and none is reduced to fuel \citep{eglash2016}. This condition has a specifically aesthetic violation that the paper will pursue under the name of the spectacle and the theory of justice (\xref{sec:distribution}), but its structure can be fixed now. There is an aesthetic extraction in which one person's perceptual labour sustains another's aesthetic consumption without return: the arranged and maintained beauty of a life or a place that is drawn upon by those who did not make it and returned to those who did as nothing, the aestheticised regard that consumes a person as scenery and gives back no share in the seeing. Such extraction is the bad cycle in the register of perception, an appropriation of aesthetic value in which the generativity of one is expropriated as the consumption of another. The good cycle is its contrary: a circulation of everyday aesthetic value in which those who generate the perception are those who benefit from it, in which the surplus returns to the shared attending that produced it, and in which no one is rendered the unattended maintainer of a beauty that others consume. To cultivate everyday aesthetic perception is therefore never a merely private refinement. It is to sustain a circulation of aesthetic value whose justice is measured by whether the value returns to all who generate it, and this is the point at which the aesthetics of the everyday becomes answerable, at its root and not by later addition, to the theory of justice.
